\documentclass[12pt, a4paper]{article}
\usepackage[utf8]{inputenc}
\usepackage[T1]{fontenc}
\usepackage[portuguese]{babel} % Adaptado para português
\usepackage{geometry}
\usepackage{graphicx}
\usepackage{tabularx}
\usepackage{booktabs}
\usepackage{float}
\usepackage{parskip}
\usepackage{enumitem}

% Configuração das margens
\geometry{
	top=2.5cm,
	bottom=2.5cm,
	left=2.5cm,
	right=2.5cm
}

\begin{document}
	
	% -------------------------------------------------------------------
	% CAPA (TITLE PAGE)
	% -------------------------------------------------------------------
	\begin{titlepage}
		\centering
		\vspace*{2cm}
		
		{\LARGE \textbf{RNG por Seeds}}\\[0.5cm]
		{\Huge \textbf{ESPECIFICAÇÃO}}\\[1.5cm]
		
		{\Large Projeto do Bloco Digital}\\[2cm]
		
		{\Large \textbf{MATHEUS GROSSI}}\\[3cm]
		
		\vfill
		
		Data: 20/01/2026\\
		Rev.: 1.0
		
		\vspace*{1cm}
	\end{titlepage}
	
	% -------------------------------------------------------------------
	% VERSÕES DO DOCUMENTO
	% -------------------------------------------------------------------
	\section*{Histórico de Versões}
	
	A tabela abaixo registra o histórico de alterações deste documento:
	
	\begin{table}[h]
		\centering
		\begin{tabularx}{\textwidth}{|c|c|X|X|}
			\hline
			\textbf{Versão} & \textbf{Data} & \textbf{Autor} & \textbf{Descrição das Mudanças} \\
			\hline
			1.0 & 20/01/2026 & Matheus Grossi & Versão inicial. Especificação completa do bloco RNG. \\
			\hline
		\end{tabularx}
	\end{table}
	
	% -------------------------------------------------------------------
	% ACRÔNIMOS
	% -------------------------------------------------------------------
	\section*{Siglas e Acrônimos}
	
	\begin{tabular}{ll}
		\textbf{RNG}  & Random Number Generator (Gerador de Números Aleatórios) \\
		\textbf{ASIC} & Application Specific Integrated Circuit \\
		\textbf{FPGA} & Field-Programmable Gate Array \\
		\textbf{FSM}  & Finite State Machine (Máquina de Estados Finita) \\
		\textbf{HDL}  & Hardware Description Language \\
		\textbf{RTL}  & Register Transfer Level \\
	\end{tabular}
	
	\newpage
	
	% -------------------------------------------------------------------
	% CONTEÚDO PRINCIPAL
	% -------------------------------------------------------------------
	
	\section{Visão Geral e Objetivos}
	
	O objetivo deste bloco digital é implementar um gerador de números de 8 bits baseado em \textit{seeds} (conjuntos pré-definidos de valores), com seleção entre quatro seeds e emissão de um novo valor sob demanda. O projeto deve assegurar:
	
	\begin{itemize}
		\item Geração de valores de 8 bits a partir de tabelas de dados associadas às seeds;
		\item Identificação da seed embutida no próprio número gerado (2 bits reservados para o identificador);
		\item Operação síncrona ao clock do sistema, com reset assíncrono ativo em nível baixo;
		\item Robustez: operação estável em condição ociosa (sem requisições) e durante reinicialização por reset;
		\item Modularidade: possibilidade de ajuste de conjuntos de valores (seeds) e políticas de emissão por meio de configuração do projeto.
	\end{itemize}
	
	\section{Identificação do Projeto (RTL Design)}
	Este design é implementado utilizando HDL padrão de indústria, destinado a rodar tanto em hardware de prototipagem quanto como um módulo em um ASIC.
	
	\begin{itemize}
		\item \textbf{Nome do bloco:} RNG por Seeds (Gerador determinístico de 8 bits)
		\item \textbf{Linguagem de Descrição:} Verilog
		\item \textbf{Tecnologias Alvo:} FPGA (Xilinx/Intel) e ASIC (Fluxo Cadence)
		\item \textbf{Data de início:} 20/01/2026
	\end{itemize}
	
	\section{Design da FSM}
	O bloco deve ser controlado por uma Máquina de Estados Finita para coordenar a emissão de valores sob requisição. Em alto nível, os estados mínimos esperados são:
	
	\begin{itemize}
		\item \textbf{IDLE:} Estado de espera/reset. Nenhuma emissão ocorre sem requisição.
		\item \textbf{REQ\_CAPTURE:} Captura/valida a requisição de novo número.
		\item \textbf{GEN:} Seleciona a seed vigente e obtém o próximo valor associado.
		\item \textbf{OUT\_UPDATE:} Atualiza a saída com o número gerado.
	\end{itemize}
	
	\section{Parâmetros e Restrições}
	
	\subsection{Requisitos de Temporização}
	O design deve fechar temporização (\textit{timing closure}) na seguinte faixa operacional:
	\begin{itemize}
		\item \textbf{Faixa:} 25MHz até 100MHz.
	\end{itemize}
	
	\subsection{Parâmetros de Configuração (Generics)}
	O projeto deve permitir configuração dos seguintes atributos:
	
	\begin{table}[H]
		\centering
		\begin{tabularx}{\textwidth}{l l l X}
			\toprule
			\textbf{Parâmetro} & \textbf{Tipo} & \textbf{Padrão} & \textbf{Descrição} \\
			\midrule
			OUT\_WIDTH & integer & 8 & Largura do número gerado (em bits). \\
			SEED\_COUNT & integer & 4 & Quantidade de seeds disponíveis. \\
			VALUES\_PER\_SEED & integer & 64 & Valores disponíveis por seed. \\
			ENABLE\_NO\_REPEAT & boolean & 0 & Habilita mitigação de repetição. \\
			\bottomrule
		\end{tabularx}
	\end{table}
	
	\section{Visão Geral do Sistema}
	
	\subsection{Interfaces do Bloco}
	A tabela abaixo descreve os sinais de entrada e saída mandatórios para o funcionamento do bloco.
	
	% FIGURA - Substitua 'example-image-a' pelo seu diagrama de blocos real se tiver
	\begin{figure}[H]
		\centering
		\includegraphics[width=0.6\textwidth]{example-image-a} 
		\caption{Diagrama de Blocos do RNG}
		\label{fig:rng_block}
	\end{figure}
	
	The following table describes inputs and outputs:
	
	\begin{table}[H]
		\centering
		\begin{tabularx}{\textwidth}{l c c X}
			\toprule
			\textbf{Nome} & \textbf{Dir} & \textbf{Bits} & \textbf{Descrição Funcional} \\
			\midrule
			clk\_i & I & 1 & Clock do sistema (borda de subida). \\
			rst\_i & I & 1 & Reset assíncrono (Active Low). \\
			start\_i & I & 1 & Habilita o fluxo de geração (início de sessão). \\
			req\_num\_i & I & 1 & Pulso de requisição de novo número. \\
			num\_to\_send\_o & O & 8 & Número gerado (valor + ID da seed). \\
			\bottomrule
		\end{tabularx}
		\caption{Definição de Interfaces (I/O)}
		\label{tab:io_desc}
	\end{table}
	
	\subsection{Lógica Interna e Sub-blocos}
	O design deve ser segregado nos seguintes módulos lógicos:
	
	\begin{itemize}
		\item \textbf{Control Unit:} coordena o atendimento de \texttt{req\_num\_i} e a atualização da saída.
		\item \textbf{Datapath:} seleciona o valor associado à seed e forma o byte de saída.
		\item \textbf{Seed Selector:} define a seed ativa (4 opções).
		\item \textbf{Mitigação (Opcional):} bloco para reduzir valores repetidos em uma sessão, controlado pelo parâmetro \texttt{ENABLE\_NO\_REPEAT}.
	\end{itemize}
	
	\subsection{Tratamento de Exceções}
	Caso ocorra um reset durante uma operação, o bloco deve descartar a transação em andamento e retornar imediatamente ao estado \textbf{IDLE}, evitando mudanças espúrias na saída durante a reinicialização.
	
	\section{Critérios de Aceite}
	O projeto será considerado concluído quando:
	\begin{enumerate}
		\item O RTL passar na verificação de sintaxe e \textit{linting} sem erros críticos.
		\item O bloco for simulável e demonstrar emissão correta de valores sob \texttt{req\_num\_i}, respeitando reset e operação ociosa.
		\item O bloco for sintetizável na tecnologia alvo sem violações de \textit{setup} ou \textit{hold} na faixa de clock especificada.
	\end{enumerate}
	
	\vfill
	\begin{center}
		\textbf{Metodologia de Design:} Synchronous Design / Boas práticas de reset e integração
	\end{center}
	
\end{document}